\chapter{CHAPTER 4: SYSTEM DESIGN AND
REQUIREMENTS}\label{chapter-4-system-design-and-requirements}

\section{4.1 Chapter Overview}\label{chapter-overview}

This chapter transitions from the methodological framework to the
explicit technical requirements and architectural design of the
SentinelAQ system. In accordance with the professor's feedback to
explicitly define the primary research inputs, Section 4.2 details the
datasets and data libraries utilized. Section 4.3 outlines the
Functional Requirements (FR) necessary for the system to achieve its
predictive and alerting goals. Section 4.4 defines the Non-Functional
Requirements (NFR) ensuring performance, scalability, and usability.
Finally, Section 4.5 presents a high-level overview of the decoupled
system architecture.

\section{4.2 Primary Research Components: Datasets and
Libraries}\label{primary-research-components-datasets-and-libraries}

As established in the methodology, this study relies on empirical
primary research rather than qualitative interviews. The foundation of
the system is built upon two programmatic data streams and a specialized
stack of Python libraries.

\subsection{4.2.1 Datasets}\label{datasets}

To construct a robust multi-modal dataset spanning from January 2022 to
December 2025, two distinct Application Programming Interfaces (APIs)
were integrated: * \textbf{PurpleAir REST API:} Provided the primary
dependent variable (ground truth PM2.5) as well as localized ambient
temperature and humidity. Data was extracted from two strategically
chosen nodes: Colombo (Coastal, ID: 29677) and Kandy (Valley, ID:
12451). * \textbf{Open-Meteo API:} Served as the critical meteorological
proxy. To mitigate the severe cloud occlusion risks associated with
satellite data (detailed in Section 1.6), Open-Meteo provided highly
reliable, hourly wind speed, wind direction, and precipitation data,
forming the backbone of the spatial analysis.

\subsection{4.2.2 Data Libraries}\label{data-libraries}

The backend infrastructure and predictive modeling were developed using
a specialized Python 3.11+ stack: * \textbf{Data Engineering:}
\texttt{pandas} and \texttt{numpy} were utilized for complex matrix
manipulations, handling missing values, and executing cyclical
trigonometric time encodings (sine/cosine transformations for hours and
months). * \textbf{Machine Learning:} \texttt{xgboost} (for tree-based
ensemble modeling), \texttt{tensorflow} and \texttt{keras} (for building
the deep learning LSTM architecture), and \texttt{statsmodels} (for the
SARIMAX statistical baseline). * \textbf{Explainable AI:} The
\texttt{shap} library was integrated specifically to compute Shapley
values, mathematically isolating the exact contribution of features like
wind speed or historical lag on individual predictions. * \textbf{Cloud
Integration:} \texttt{firebase-admin} facilitated secure backend
communication with Google Cloud, handling both Firestore database writes
and Cloud Messaging (FCM) notifications.

\section{4.3 Functional Requirements
(FR)}\label{functional-requirements-fr}

Functional requirements define the specific, testable behaviors the
SentinelAQ system must execute to deliver end-to-end value. *
\textbf{FR1 (Automated Data Ingestion):} The backend Python pipeline
shall automatically fetch the latest hourly PM2.5 and meteorological
readings from the PurpleAir and Open-Meteo APIs without manual
intervention. * \textbf{FR2 (Dynamic Feature Engineering):} The system
shall dynamically construct the 48-dimensional feature vector in
real-time, calculating 3-hour, 6-hour, and 24-hour rolling statistical
means from the newly ingested data. * \textbf{FR3 (Multi-Horizon
Inference):} The system shall execute the pre-trained XGBoost model to
generate distinct PM2.5 concentration forecasts for 1, 6, 12, 24, and
48-hour future horizons. * \textbf{FR4 (Real-Time Explainability):} The
system shall compute SHAP \texttt{TreeExplainer} values for every
forecast, explicitly isolating the top three atmospheric drivers (e.g.,
``Driven by: 24h Accumulation \& Low Wind Speed''). * \textbf{FR5
(Database Synchronization):} The backend shall structure the predictions
and SHAP interpretations into JSON payloads and execute synchronous
write operations to the Firebase Firestore cloud database. * \textbf{FR6
(Mobile Visualization):} The React Native frontend shall fetch the
latest Firestore document and render the current AQI via a dynamic
color-coded dial, alongside a responsive 48-hour trend chart. *
\textbf{FR7 (Threshold Alerting):} The system shall allow users to
configure personalized AQI thresholds (e.g., ``Alert me if PM2.5
\textgreater{} 50''). The backend shall trigger Firebase Cloud Messaging
(FCM) push notifications when forecasted levels breach these
user-defined thresholds.

\section{4.4 Non-Functional Requirements
(NFR)}\label{non-functional-requirements-nfr}

Non-functional requirements specify the critical quality attributes and
constraints under which the system must operate. * \textbf{NFR1
(Implementation Speed \& Performance):} The automated backend inference
pipeline must complete the entire cycle---data extraction, vector
engineering, SHAP calculation, and database synchronization---in under
60 seconds to ensure the mobile client receives truly real-time updates.
* \textbf{NFR2 (High Availability):} Leveraging Google Firebase's
serverless infrastructure, the system database and messaging services
must maintain 99.9\% uptime, ensuring users have uninterrupted access to
health advisories during severe pollution events. * \textbf{NFR3
(Architectural Scalability):} The Firestore NoSQL database must be
structurally optimized (document-collection architecture) to easily
accommodate the addition of new geographical nodes (e.g., expanding to
Jaffna or Galle) without requiring a complete schema redesign or causing
data retrieval bottlenecks. * \textbf{NFR4 (Usability \& Cognitive
Load):} The mobile application must instantly translate complex
microgram (µg/m³) predictions into intuitive, universally understood
visual indicators (US EPA AQI standard colors: Green, Yellow, Orange,
Red) to minimize cognitive load for non-technical users.

\section{4.5 System Architecture
Overview}\label{system-architecture-overview}

To satisfy the aforementioned requirements, SentinelAQ was designed
utilizing a decoupled, four-lane enterprise architecture: 1.
\textbf{User Lane:} Represents the physical actor configuring alert
thresholds and consuming visualizations. 2. \textbf{Mobile Client
(Frontend):} A React Native application responsible for handling local
UI states, rendering the AQI dial, and displaying push notifications. 3.
\textbf{API Backend (Middleware):} The Python-driven core where the
XGBoost inference and SHAP explainability matrices are processed. 4.
\textbf{External Providers (Cloud/Data):} The integration layer bridging
third-party sensor networks (PurpleAir), weather APIs (Open-Meteo), and
cloud hosting (Firebase).

This decoupled design ensures that the heavy computational load of
machine learning inference is entirely abstracted away from the user's
mobile device, allowing the app to remain lightweight and highly
responsive. The following chapter (Chapter 5) details the iterative
journey of developing this architecture, documenting the failures and
successes encountered across three distinct versions (v1 to v3).
